
\documentclass[a4paper,10pt]{article}
\usepackage[margin=7mm]{geometry}
\usepackage{fontspec}
\usepackage{xeCJK}
\usepackage{graphicx}
\usepackage{multicol}
\usepackage{amsmath,amssymb}
\usepackage{enumitem}
\usepackage{hyperref}
\usepackage{xcolor}
\usepackage{titlesec}
\setmainfont{Noto Sans}
\setsansfont{Noto Sans}
\setCJKmainfont{Noto Sans CJK JP}
\definecolor{TUBSred}{HTML}{BE1E3C}
\setlength{\parindent}{0pt}
\setlength{\parskip}{2.8pt}
\setlength{\columnsep}{5mm}
\setlist[itemize]{leftmargin=1.25em,itemsep=1.5pt,topsep=2pt,parsep=0pt}
\setlist[enumerate]{leftmargin=1.45em,itemsep=2pt,topsep=2pt,parsep=0pt}
\hypersetup{colorlinks=true,linkcolor=blue,urlcolor=blue}
\titleformat{\section}{\Large\bfseries\color{TUBSred}}{}{0pt}{}
\newcommand{\SlideHeader}[3]{%
  \noindent{\large\bfseries #1 / Slide #2: #3}\hfill {\scriptsize review v2}\par
  \vspace{1mm}\hrule\vspace{1.5mm}
}
\newcommand{\SlideImage}[2]{%
  \noindent\makebox[\textwidth][c]{\includegraphics[page=#1,width=0.96\textwidth,height=0.36\textheight,keepaspectratio]{#2}}\par
}
\newcommand{\NoteSection}[1]{%
  \vspace{2.2mm}\noindent{\normalsize\bfseries\color{TUBSred}#1}\par\vspace{0.8mm}
}
\newcommand{\ExamItem}[3]{%
  \vspace{1.2mm}\noindent{\scriptsize\bfseries #1}\par
  \noindent{\scriptsize\bfseries 問題内容:} #2\par
  \noindent{\scriptsize\bfseries 期待される回答:} #3\par
}
\begin{document}
\begin{titlepage}
\centering
\vspace*{2cm}
{\Huge Data Driven Decision Making\par}
\vspace{0.5cm}
{\Large Lecture 1: Overview - 日本語訳・解説・試験対策ノート\par}
\vspace{1cm}
\begin{minipage}{0.88\textwidth}
このPDFは、元スライドを保持しつつ、日本語訳、概念解説、試験対策ポイント、過去問関連、確認問題と解答を構造化したレビュー用ノートです。日本語訳はスライド構造を保ち、解説は初学者が試験答案で使える水準を目標にしています。過去問は出題範囲を限定するものではなく、重要論点を確認するための参考です。
\end{minipage}
\vfill
{\large Materials/2026}\par
\end{titlepage}
\tableofcontents
\clearpage
\section*{Lecture 1: Overview}

\addcontentsline{toc}{section}{Lecture 1: Overview}

\clearpage

\SlideHeader{Lecture 1: Overview}{001}{表紙}

\SlideImage{1}{../Materials/2026/3DM_01_Overview_SS26.pdf}

\vspace{1.5mm}

\begin{multicols}{2}

\footnotesize

\raggedcolumns

\NoteSection{日本語訳}

\begin{itemize}
\item Data Driven Decision Making
\item 第1回 - 概要
\end{itemize}

\end{multicols}

\clearpage

\SlideHeader{Lecture 1: Overview}{002}{連絡先}

\SlideImage{2}{../Materials/2026/3DM_01_Overview_SS26.pdf}

\vspace{1.5mm}

\begin{multicols}{2}

\footnotesize

\raggedcolumns

\NoteSection{日本語訳}

\begin{itemize}
\item 担当教員
\begin{itemize}
\item Prof. Dr. Dirk Mattfeld
\item オフィスアワー: 登録制
\item d.mattfeld@tu-braunschweig.de
\end{itemize}
\item ティーチングアシスタント
\begin{itemize}
\item Felix Moeller
\item オフィスアワー: 登録制
\item felix.moeller@tu-braunschweig.de
\end{itemize}
\item Decision Support オフィス
\begin{itemize}
\item オフィスアワー: Webサイトで告知
\item 電話: (0531) 391-3211
\item ds@tu-braunschweig.de
\end{itemize}
\item Webサイト
\begin{itemize}
\item 講義・試験日程などを確認する場所
\end{itemize}
\end{itemize}

\end{multicols}

\clearpage

\SlideHeader{Lecture 1: Overview}{003}{Decision Support の修士科目}

\SlideImage{3}{../Materials/2026/3DM_01_Overview_SS26.pdf}

\vspace{1.5mm}

\begin{multicols}{2}

\footnotesize

\raggedcolumns

\NoteSection{日本語訳}

\begin{itemize}
\item 基礎・方向づけ
\begin{itemize}
\item Decision Support (5 LP)
\item Intelligent Data Analysis (WS)
\item Planning for Mobility and Transportation Purposes (WS)
\end{itemize}
\item 専門
\begin{itemize}
\item Decision Support (5 LP)
\item Data Driven Decision Making (SS)
\item Exercise in Data Driven Decision Making (SS)
\item Master seminar
\item Master thesis (6 months)
\end{itemize}
\item その他の専攻、とくに機械工学系学生にも開かれている。
\end{itemize}

\end{multicols}

\clearpage

\SlideHeader{Lecture 1: Overview}{004}{講義について}

\SlideImage{4}{../Materials/2026/3DM_01_Overview_SS26.pdf}

\vspace{1.5mm}

\begin{multicols}{2}

\footnotesize

\raggedcolumns

\NoteSection{日本語訳}

\begin{itemize}
\item Data Driven Decision Making
\item 以前は『モビリティ応用のための情報システム』、さらに以前は『交通情報システム』だった。
\item 焦点は、情報システムの描写から、情報システム内の自動意思決定へ移っている。
\item Uber, Lime, Amazon などの先端的デジタルビジネスモデルと整合する。
\item PMT は意思決定に、IDA はデータ分析に焦点を置く。DDDM はその上に構築される。
\item 3DM は演習とともに『Decision Support専門』を構成する。
\end{itemize}

\NoteSection{解説}
DDDM は、データを単に可視化・分析するだけでなく、そのデータを使って具体的な行動を選ぶための枠組みである。したがって中心は『どの情報を使い、どのモデルで、どの決定を選び、現実にどう実装するか』である。\par

stochasticity は、需要、交通、サービス時間などの情報が不確実で、時間とともに新しく明らかになることを指す。dynamism は、新しい情報が出た後に計画を更新しながら意思決定を続ける性質を指す。\par

\NoteSection{試験対策ポイント}
DDDMは、データ分析、意思決定モデル、現実への実装を結ぶ講義である。stochasticity と dynamism を必ず区別して説明できるようにする。\par

\NoteSection{確認問題と解答}
\begin{enumerate}
\item \textbf{DDDM は通常のデータ分析と何が違うか。}\\
通常のデータ分析は、過去データの説明や予測で終わることがある。DDDMでは、その情報を使って実際の決定、例えば配送順序、車両割当、注文受諾、在庫補充などを選ぶ。したがって、情報モデルと意思決定モデル、さらに現実への実装までを一体として考える。
\item \textbf{stochasticity と dynamism を区別して説明せよ。}\\
stochasticity は将来情報や実現値が不確実であること、dynamism は時間の経過とともに状態が変わり、そのたびに決定を更新する必要があることである。例えば、注文が確率的に発生することは stochasticity、注文発生後に配送計画を組み替えることは dynamism である。
\end{enumerate}

\end{multicols}

\clearpage

\SlideHeader{Lecture 1: Overview}{005}{目標}

\SlideImage{5}{../Materials/2026/3DM_01_Overview_SS26.pdf}

\vspace{1.5mm}

\begin{multicols}{2}

\footnotesize

\raggedcolumns

\NoteSection{日本語訳}

\begin{itemize}
\item 確率性と動的性質を理解する。
\item 確率的・動的意思決定問題をモデル化する。
\item 方法論の概要、機能、強み、弱みを理解する。
\item デジタル経済における関連問題を理解する。
\item 60分筆記試験に合格する。
\item 修士論文、博士課程への準備を行う。
\end{itemize}

\NoteSection{解説}
DDDM は、データを単に可視化・分析するだけでなく、そのデータを使って具体的な行動を選ぶための枠組みである。したがって中心は『どの情報を使い、どのモデルで、どの決定を選び、現実にどう実装するか』である。\par

stochasticity は、需要、交通、サービス時間などの情報が不確実で、時間とともに新しく明らかになることを指す。dynamism は、新しい情報が出た後に計画を更新しながら意思決定を続ける性質を指す。\par

\NoteSection{試験対策ポイント}
DDDMは、データ分析、意思決定モデル、現実への実装を結ぶ講義である。stochasticity と dynamism を必ず区別して説明できるようにする。\par

\NoteSection{確認問題と解答}
\begin{enumerate}
\item \textbf{DDDM は通常のデータ分析と何が違うか。}\\
通常のデータ分析は、過去データの説明や予測で終わることがある。DDDMでは、その情報を使って実際の決定、例えば配送順序、車両割当、注文受諾、在庫補充などを選ぶ。したがって、情報モデルと意思決定モデル、さらに現実への実装までを一体として考える。
\item \textbf{stochasticity と dynamism を区別して説明せよ。}\\
stochasticity は将来情報や実現値が不確実であること、dynamism は時間の経過とともに状態が変わり、そのたびに決定を更新する必要があることである。例えば、注文が確率的に発生することは stochasticity、注文発生後に配送計画を組み替えることは dynamism である。
\end{enumerate}

\end{multicols}

\clearpage

\SlideHeader{Lecture 1: Overview}{006}{DDDMの講義計画}

\SlideImage{6}{../Materials/2026/3DM_01_Overview_SS26.pdf}

\vspace{1.5mm}

\begin{multicols}{2}

\footnotesize

\raggedcolumns

\NoteSection{日本語訳}

\begin{itemize}
\item 講義トピック
\begin{itemize}
\item 1 概要
\item 2 情報モデリングと意思決定モデリング
\item 3 確率性
\item 4 動的性
\item 5 マルコフ意思決定過程
\item 6 近似動的計画法
\item 7 Look-ahead Methods
\item 8 価値関数近似
\item 9 ADPにおける分析上の問題
\item 10 Combined Methods
\item 11 プラットフォーム型フードデリバリー
\end{itemize}
\end{itemize}

\end{multicols}

\clearpage

\SlideHeader{Lecture 1: Overview}{007}{文献}

\SlideImage{7}{../Materials/2026/3DM_01_Overview_SS26.pdf}

\vspace{1.5mm}

\begin{multicols}{2}

\footnotesize

\raggedcolumns

\NoteSection{日本語訳}

\begin{itemize}
\item Warren Powell の研究。
\item TU Braunschweig における Marlin Ulmer らとの近年の共同研究。
\item 交通・モビリティに関する適用と発展。
\item Ulmer et al. (2020): stochastic dynamic vehicle routing problems のモデル化。
\item Ulmer et al. (2021): prescriptive analytics の観点から見た stochastic dynamic vehicle routing のレビュー。
\end{itemize}

\end{multicols}

\clearpage

\SlideHeader{Lecture 1: Overview}{008}{演習について}

\SlideImage{8}{../Materials/2026/3DM_01_Overview_SS26.pdf}

\vspace{1.5mm}

\begin{multicols}{2}

\footnotesize

\raggedcolumns

\NoteSection{日本語訳}

\begin{itemize}
\item 『Decision Support専門』の Studienleistung (2.5 LP)。
\item serious game アプリケーション『Trucks and Barges』で講義内容を練習する。
\item シミュレーション内容
\begin{itemize}
\item 貨物業者が、複数日にわたり、ターミナルからコンテナを異なる輸送モードで出荷する意思決定状況。
\end{itemize}
\item 運営
\begin{itemize}
\item 内容が積み上がる5回のオンライン日程。
\item 宿題。
\item StudIP: Data Driven Decision Making - Uebung
\end{itemize}
\end{itemize}

\end{multicols}

\clearpage

\SlideHeader{Lecture 1: Overview}{009}{試験情報}

\SlideImage{9}{../Materials/2026/3DM_01_Overview_SS26.pdf}

\vspace{1.5mm}

\begin{multicols}{2}

\footnotesize

\raggedcolumns

\NoteSection{日本語訳}

\begin{itemize}
\item 試験会場と時刻は講座のWebサイトで確認できる。
\item 情報は定期的に更新される。
\item 試験日は早めに告知され、試験教室は試験直前に割り当てられる。
\item 必要な情報がWebサイトにない場合のみメールで問い合わせる。
\end{itemize}

\end{multicols}

\clearpage

\SlideHeader{Lecture 1: Overview}{010}{過去問}

\SlideImage{10}{../Materials/2026/3DM_01_Overview_SS26.pdf}

\vspace{1.5mm}

\begin{multicols}{2}

\footnotesize

\raggedcolumns

\NoteSection{日本語訳}

\begin{itemize}
\item 近年の過去問の印刷版は研究所で入手できる。部数には限りがある。
\item 過去問はWebサイトでも確認できる。
\item 過去問を解くだけでは十分な準備ではなく、良い成績を保証しない。
\item 講義で扱った全トピックが試験対象である。
\item 問題形式は似ることがあるが、内容は変わり、過去問にない新問も含まれ得る。
\item 過去問の模範解答は提供されない。
\end{itemize}

\end{multicols}

\clearpage

\SlideHeader{Lecture 1: Overview}{011}{アジェンダ}

\SlideImage{11}{../Materials/2026/3DM_01_Overview_SS26.pdf}

\vspace{1.5mm}

\begin{multicols}{2}

\footnotesize

\raggedcolumns

\NoteSection{日本語訳}

\begin{itemize}
\item 動機
\item Operations Research と Business Analytics の復習
\item モデリングと講義全体の見取り図
\end{itemize}

\end{multicols}

\clearpage

\SlideHeader{Lecture 1: Overview}{012}{動機}

\SlideImage{12}{../Materials/2026/3DM_01_Overview_SS26.pdf}

\vspace{1.5mm}

\begin{multicols}{2}

\footnotesize

\raggedcolumns

\NoteSection{日本語訳}

\begin{itemize}
\item 新しく変化するビジネスコンセプト: より速く、より個別的に、リアルタイムに。
\item 例
\begin{itemize}
\item 即時配送
\item オンデマンド生産・受注生産
\item 予約スケジューリング
\item その他
\end{itemize}
\item 計画と制御が難しくなる。
\item 二つの共通テーマ
\begin{itemize}
\item Stochasticity: 情報が時間とともに変化する。
\item Dynamism: 情報変化を踏まえて計画を更新する。
\end{itemize}
\item 確率的・動的意思決定過程へつながる。
\end{itemize}

\NoteSection{解説}
DDDM は、データを単に可視化・分析するだけでなく、そのデータを使って具体的な行動を選ぶための枠組みである。したがって中心は『どの情報を使い、どのモデルで、どの決定を選び、現実にどう実装するか』である。\par

stochasticity は、需要、交通、サービス時間などの情報が不確実で、時間とともに新しく明らかになることを指す。dynamism は、新しい情報が出た後に計画を更新しながら意思決定を続ける性質を指す。\par

\NoteSection{試験対策ポイント}
DDDMは、データ分析、意思決定モデル、現実への実装を結ぶ講義である。stochasticity と dynamism を必ず区別して説明できるようにする。\par

\NoteSection{確認問題と解答}
\begin{enumerate}
\item \textbf{DDDM は通常のデータ分析と何が違うか。}\\
通常のデータ分析は、過去データの説明や予測で終わることがある。DDDMでは、その情報を使って実際の決定、例えば配送順序、車両割当、注文受諾、在庫補充などを選ぶ。したがって、情報モデルと意思決定モデル、さらに現実への実装までを一体として考える。
\item \textbf{stochasticity と dynamism を区別して説明せよ。}\\
stochasticity は将来情報や実現値が不確実であること、dynamism は時間の経過とともに状態が変わり、そのたびに決定を更新する必要があることである。例えば、注文が確率的に発生することは stochasticity、注文発生後に配送計画を組み替えることは dynamism である。
\end{enumerate}

\end{multicols}

\clearpage

\SlideHeader{Lecture 1: Overview}{013}{例: 輸送}

\SlideImage{13}{../Materials/2026/3DM_01_Overview_SS26.pdf}

\vspace{1.5mm}

\begin{multicols}{2}

\footnotesize

\raggedcolumns

\NoteSection{日本語訳}

\begin{itemize}
\item 同日配送、ピックアップ、ドローン配送、食品配送など、輸送サービスの例。
\end{itemize}

\NoteSection{解説}
DDDM は、データを単に可視化・分析するだけでなく、そのデータを使って具体的な行動を選ぶための枠組みである。したがって中心は『どの情報を使い、どのモデルで、どの決定を選び、現実にどう実装するか』である。\par

stochasticity は、需要、交通、サービス時間などの情報が不確実で、時間とともに新しく明らかになることを指す。dynamism は、新しい情報が出た後に計画を更新しながら意思決定を続ける性質を指す。\par

\NoteSection{試験対策ポイント}
DDDMは、データ分析、意思決定モデル、現実への実装を結ぶ講義である。stochasticity と dynamism を必ず区別して説明できるようにする。\par

\NoteSection{確認問題と解答}
\begin{enumerate}
\item \textbf{DDDM は通常のデータ分析と何が違うか。}\\
通常のデータ分析は、過去データの説明や予測で終わることがある。DDDMでは、その情報を使って実際の決定、例えば配送順序、車両割当、注文受諾、在庫補充などを選ぶ。したがって、情報モデルと意思決定モデル、さらに現実への実装までを一体として考える。
\item \textbf{stochasticity と dynamism を区別して説明せよ。}\\
stochasticity は将来情報や実現値が不確実であること、dynamism は時間の経過とともに状態が変わり、そのたびに決定を更新する必要があることである。例えば、注文が確率的に発生することは stochasticity、注文発生後に配送計画を組み替えることは dynamism である。
\end{enumerate}

\end{multicols}

\clearpage

\SlideHeader{Lecture 1: Overview}{014}{例: モビリティとサービス}

\SlideImage{14}{../Materials/2026/3DM_01_Overview_SS26.pdf}

\vspace{1.5mm}

\begin{multicols}{2}

\footnotesize

\raggedcolumns

\NoteSection{日本語訳}

\begin{itemize}
\item バイクシェア、オンデマンド交通、医療訪問、タクシー、エレベーター保守、緊急車両など、モビリティ・サービスの例。
\end{itemize}

\NoteSection{解説}
DDDM は、データを単に可視化・分析するだけでなく、そのデータを使って具体的な行動を選ぶための枠組みである。したがって中心は『どの情報を使い、どのモデルで、どの決定を選び、現実にどう実装するか』である。\par

stochasticity は、需要、交通、サービス時間などの情報が不確実で、時間とともに新しく明らかになることを指す。dynamism は、新しい情報が出た後に計画を更新しながら意思決定を続ける性質を指す。\par

\NoteSection{試験対策ポイント}
DDDMは、データ分析、意思決定モデル、現実への実装を結ぶ講義である。stochasticity と dynamism を必ず区別して説明できるようにする。\par

\NoteSection{確認問題と解答}
\begin{enumerate}
\item \textbf{DDDM は通常のデータ分析と何が違うか。}\\
通常のデータ分析は、過去データの説明や予測で終わることがある。DDDMでは、その情報を使って実際の決定、例えば配送順序、車両割当、注文受諾、在庫補充などを選ぶ。したがって、情報モデルと意思決定モデル、さらに現実への実装までを一体として考える。
\item \textbf{stochasticity と dynamism を区別して説明せよ。}\\
stochasticity は将来情報や実現値が不確実であること、dynamism は時間の経過とともに状態が変わり、そのたびに決定を更新する必要があることである。例えば、注文が確率的に発生することは stochasticity、注文発生後に配送計画を組み替えることは dynamism である。
\end{enumerate}

\end{multicols}

\clearpage

\SlideHeader{Lecture 1: Overview}{015}{まとめ}

\SlideImage{15}{../Materials/2026/3DM_01_Overview_SS26.pdf}

\vspace{1.5mm}

\begin{multicols}{2}

\footnotesize

\raggedcolumns

\NoteSection{日本語訳}

\begin{itemize}
\item サービスプロセスは時間を通じて進行する。
\item 需要、依頼、交通、ドライバー、バッテリー残量、要件などの新情報が発生する。
\item 計画への反応と更新が必要になる。
\item 反応が遅すぎる場合がある。例: 車両が渋滞に巻き込まれる、食事配送のドライバーが周辺にいない、ドローンが墜落する。
\item 期待される将来を先読みする必要がある。
\item 将来のシステム状態に対する柔軟性を目指す。
\end{itemize}

\NoteSection{解説}
DDDM は、データを単に可視化・分析するだけでなく、そのデータを使って具体的な行動を選ぶための枠組みである。したがって中心は『どの情報を使い、どのモデルで、どの決定を選び、現実にどう実装するか』である。\par

stochasticity は、需要、交通、サービス時間などの情報が不確実で、時間とともに新しく明らかになることを指す。dynamism は、新しい情報が出た後に計画を更新しながら意思決定を続ける性質を指す。\par

\NoteSection{試験対策ポイント}
DDDMは、データ分析、意思決定モデル、現実への実装を結ぶ講義である。stochasticity と dynamism を必ず区別して説明できるようにする。\par

\NoteSection{確認問題と解答}
\begin{enumerate}
\item \textbf{DDDM は通常のデータ分析と何が違うか。}\\
通常のデータ分析は、過去データの説明や予測で終わることがある。DDDMでは、その情報を使って実際の決定、例えば配送順序、車両割当、注文受諾、在庫補充などを選ぶ。したがって、情報モデルと意思決定モデル、さらに現実への実装までを一体として考える。
\item \textbf{stochasticity と dynamism を区別して説明せよ。}\\
stochasticity は将来情報や実現値が不確実であること、dynamism は時間の経過とともに状態が変わり、そのたびに決定を更新する必要があることである。例えば、注文が確率的に発生することは stochasticity、注文発生後に配送計画を組み替えることは dynamism である。
\end{enumerate}

\end{multicols}

\clearpage

\SlideHeader{Lecture 1: Overview}{016}{意思決定}

\SlideImage{16}{../Materials/2026/3DM_01_Overview_SS26.pdf}

\vspace{1.5mm}

\begin{multicols}{2}

\footnotesize

\raggedcolumns

\NoteSection{日本語訳}

\begin{itemize}
\item Real-World: 行動が実行され、その結果として状態を観測できる。
\item Aggregation: 現実世界のデータを情報モデルの次元へ移す。
\item Information Model: 情報を圧縮された形で描写する。
\item Decision Model: 決定空間を定義し、決定を評価する。
\item Implementation: 意思決定モデルの解を現実世界の行動へ変換する。
\end{itemize}

\NoteSection{解説}
Real-World は、配送、交通、顧客、車両、作業員など、実際に行動が行われる対象である。ここでは全ての詳細をモデルに入れることはできないため、意思決定に必要な情報だけを選ぶ。\par

Aggregation は、観測データをモデルで使える次元へ変換する処理である。例えばGPS位置列を道路セグメント別の平均速度へ変換すること、住所を顧客ノードへ変換することが該当する。\par

Information Model は、現実を圧縮した情報表現である。旅行時間行列、時間帯別速度パターン、顧客集合、需要分布などが含まれる。\par

Decision Model は、可能な決定、目的関数、制約、評価方法を定義する。TSP/VRPであれば、訪問順序、利用アーク、各顧客を一度訪問する制約、総旅行時間最小化が含まれる。\par

Implementation は、モデルの解を現実の行動へ戻す処理である。例えば『1-3-2』というノード順序を、実際の住所と道路ナビゲーションへ変換する。\par

\NoteSection{試験対策ポイント}
Real-World, Aggregation, Information Model, Decision Model, Implementation の役割とデータ・解の流れを、配送例に対応付けて説明できること。\par

\NoteSection{過去問関連}
\ExamItem{WS23/24 Aufgabe 2 (5点)}{Real-World Problem と Decision Model の関係、および Information Model を介した流れを説明する問題。}{現実世界から直接最適化モデルを作るのではなく、まず観測データを集約して意思決定に必要な情報表現に落とす。Information Model は現実を圧縮した表現で、Decision Model はその表現を使って決定空間、制約、目的関数を定義する。解は Implementation により現実の行動へ変換され、結果観測から仮説・モデルが更新される。}
\ExamItem{SS24 Aufgabe 1a (8点)}{Information Model, Decision Model, Aggregation, Disaggregation/Implementation の機能を説明する問題。}{Aggregation は現実データをモデルの次元へ変換する。Information Model は意思決定に必要な状態・属性・パラメータを整理する。Decision Model は可能な決定、目的、制約、評価方法を定義する。Disaggregation/Implementation はモデル解を現実の作業指示やルートへ戻す。重要なのは、これらが一方向で終わらず、実装後の観測で更新されるループである点である。}
\ExamItem{WS22/23 Aufgabe 3 (6点)}{Business Analytics プロセスの中で OR 的なモデル化・解法がどこにあるかを説明する問題。}{情報モデルはデータ側、意思決定モデルはOR側の橋渡しである。ORは、現実の問題構造を抽象化し、目的関数・制約・決定変数に変換し、解を評価可能にする。}
\ExamItem{WS22/23 Aufgabe 2 (6点)}{Business Analytics を構成する分野と、それらの交差領域を説明する問題。}{Operations Research は意思決定モデル、制約、目的関数、最適化を担う。Statistics はデータの不確実性、分布、推定、予測を扱う。Computer Science/Information Systems はデータ管理、アルゴリズム、実装、情報システムを担う。Business Analytics はこれらを組み合わせ、データから意思決定可能な知見と実行可能な解を作る分野である。}
\ExamItem{WS22/23 Aufgabe 3 (6点)}{Business Analytics の適用プロセスを挙げ、Operations Research に属する工程を示す問題。}{現実問題を観察し、データを取得・集約して情報モデルを作り、そこから意思決定モデルを定義し、解法で解を求め、実装して現実に戻す。ORに特に属するのは、目的関数・制約・決定変数を持つ意思決定モデルの定式化、解空間の探索、最適化・ヒューリスティックによる解の計算である。}
\ExamItem{WS23/24 Aufgabe 1 (6点)}{Business Analytics の三つの構成分野と交差領域を説明する問題。}{三分野を単に列挙するだけでなく、データ分析だけでは意思決定にならず、最適化だけでも現実データを反映できない点を述べる。交差領域では、データからモデルを作り、モデルから現実に実装可能な意思決定を作ることが中心になる。}
\ExamItem{WS23/24 Aufgabe 3 (10点)}{配送車両の業務的ツアープランニングに Business Analytics の構成要素を対応付ける問題。}{Real-world は配送現場、顧客、道路、車両である。Aggregation は住所・道路・走行時間観測をモデル変数へ変換すること。Information model は顧客集合、距離/旅行時間行列、時間帯情報である。Decision model は訪問順序・経路・目的関数・制約を持つVRP/TSPである。Implementation は得られたツアーをドライバーのルート案内として実行することである。}

\NoteSection{確認問題と解答}
\begin{enumerate}
\item \textbf{Information Model と Decision Model の違いを説明せよ。}\\
Information Model は、意思決定に必要な現実情報の表現である。例として顧客、道路、旅行時間行列、時間帯別速度がある。Decision Model は、その情報を用いて決定を選ぶための目的関数、制約、決定変数の体系である。情報モデルが『何を知っているか』を表し、意思決定モデルが『その情報で何をどう決めるか』を表す。
\item \textbf{Aggregation と Implementation の役割を説明せよ。}\\
Aggregation は現実データを抽象化してモデル入力に変換する処理であり、Implementation はモデル解を実際の行動へ具体化する処理である。前者は現実からモデルへの橋、後者はモデルから現実への橋である。
\item \textbf{なぜこのループは一回で終わらないのか。}\\
実装後に現実の結果が観測されるためである。予測より配送が遅れた、交通パターンが違った、制約が現実に合わなかった場合、仮説、情報モデル、意思決定モデルを更新する必要がある。
\end{enumerate}

\end{multicols}

\clearpage

\SlideHeader{Lecture 1: Overview}{017}{DDDMにおける意思決定}

\SlideImage{17}{../Materials/2026/3DM_01_Overview_SS26.pdf}

\vspace{1.5mm}

\begin{multicols}{2}

\footnotesize

\raggedcolumns

\NoteSection{日本語訳}

\begin{itemize}
\item DDDMでは、情報モデルだけを詳しく扱うわけではない。
\item 最適化モデルだけに焦点を当てるわけでもない。
\item 現実世界のプロセスを含むループ全体を考慮する。
\item そのため、システムを時間を通じてモデル化する必要がある。これは dynamic である。
\item 現実世界は外部効果により変化し得る。これは stochastic である。
\item 実装された決定は、すでに変化した現実世界と衝突する可能性がある。
\end{itemize}

\NoteSection{解説}
Real-World は、配送、交通、顧客、車両、作業員など、実際に行動が行われる対象である。ここでは全ての詳細をモデルに入れることはできないため、意思決定に必要な情報だけを選ぶ。\par

Aggregation は、観測データをモデルで使える次元へ変換する処理である。例えばGPS位置列を道路セグメント別の平均速度へ変換すること、住所を顧客ノードへ変換することが該当する。\par

Information Model は、現実を圧縮した情報表現である。旅行時間行列、時間帯別速度パターン、顧客集合、需要分布などが含まれる。\par

Decision Model は、可能な決定、目的関数、制約、評価方法を定義する。TSP/VRPであれば、訪問順序、利用アーク、各顧客を一度訪問する制約、総旅行時間最小化が含まれる。\par

Implementation は、モデルの解を現実の行動へ戻す処理である。例えば『1-3-2』というノード順序を、実際の住所と道路ナビゲーションへ変換する。\par

\NoteSection{試験対策ポイント}
Real-World, Aggregation, Information Model, Decision Model, Implementation の役割とデータ・解の流れを、配送例に対応付けて説明できること。\par

\NoteSection{過去問関連}
\ExamItem{WS23/24 Aufgabe 2 (5点)}{Real-World Problem と Decision Model の関係、および Information Model を介した流れを説明する問題。}{現実世界から直接最適化モデルを作るのではなく、まず観測データを集約して意思決定に必要な情報表現に落とす。Information Model は現実を圧縮した表現で、Decision Model はその表現を使って決定空間、制約、目的関数を定義する。解は Implementation により現実の行動へ変換され、結果観測から仮説・モデルが更新される。}
\ExamItem{SS24 Aufgabe 1a (8点)}{Information Model, Decision Model, Aggregation, Disaggregation/Implementation の機能を説明する問題。}{Aggregation は現実データをモデルの次元へ変換する。Information Model は意思決定に必要な状態・属性・パラメータを整理する。Decision Model は可能な決定、目的、制約、評価方法を定義する。Disaggregation/Implementation はモデル解を現実の作業指示やルートへ戻す。重要なのは、これらが一方向で終わらず、実装後の観測で更新されるループである点である。}
\ExamItem{WS22/23 Aufgabe 3 (6点)}{Business Analytics プロセスの中で OR 的なモデル化・解法がどこにあるかを説明する問題。}{情報モデルはデータ側、意思決定モデルはOR側の橋渡しである。ORは、現実の問題構造を抽象化し、目的関数・制約・決定変数に変換し、解を評価可能にする。}
\ExamItem{WS22/23 Aufgabe 2 (6点)}{Business Analytics を構成する分野と、それらの交差領域を説明する問題。}{Operations Research は意思決定モデル、制約、目的関数、最適化を担う。Statistics はデータの不確実性、分布、推定、予測を扱う。Computer Science/Information Systems はデータ管理、アルゴリズム、実装、情報システムを担う。Business Analytics はこれらを組み合わせ、データから意思決定可能な知見と実行可能な解を作る分野である。}
\ExamItem{WS22/23 Aufgabe 3 (6点)}{Business Analytics の適用プロセスを挙げ、Operations Research に属する工程を示す問題。}{現実問題を観察し、データを取得・集約して情報モデルを作り、そこから意思決定モデルを定義し、解法で解を求め、実装して現実に戻す。ORに特に属するのは、目的関数・制約・決定変数を持つ意思決定モデルの定式化、解空間の探索、最適化・ヒューリスティックによる解の計算である。}
\ExamItem{WS23/24 Aufgabe 1 (6点)}{Business Analytics の三つの構成分野と交差領域を説明する問題。}{三分野を単に列挙するだけでなく、データ分析だけでは意思決定にならず、最適化だけでも現実データを反映できない点を述べる。交差領域では、データからモデルを作り、モデルから現実に実装可能な意思決定を作ることが中心になる。}
\ExamItem{WS23/24 Aufgabe 3 (10点)}{配送車両の業務的ツアープランニングに Business Analytics の構成要素を対応付ける問題。}{Real-world は配送現場、顧客、道路、車両である。Aggregation は住所・道路・走行時間観測をモデル変数へ変換すること。Information model は顧客集合、距離/旅行時間行列、時間帯情報である。Decision model は訪問順序・経路・目的関数・制約を持つVRP/TSPである。Implementation は得られたツアーをドライバーのルート案内として実行することである。}

\NoteSection{確認問題と解答}
\begin{enumerate}
\item \textbf{Information Model と Decision Model の違いを説明せよ。}\\
Information Model は、意思決定に必要な現実情報の表現である。例として顧客、道路、旅行時間行列、時間帯別速度がある。Decision Model は、その情報を用いて決定を選ぶための目的関数、制約、決定変数の体系である。情報モデルが『何を知っているか』を表し、意思決定モデルが『その情報で何をどう決めるか』を表す。
\item \textbf{Aggregation と Implementation の役割を説明せよ。}\\
Aggregation は現実データを抽象化してモデル入力に変換する処理であり、Implementation はモデル解を実際の行動へ具体化する処理である。前者は現実からモデルへの橋、後者はモデルから現実への橋である。
\item \textbf{なぜこのループは一回で終わらないのか。}\\
実装後に現実の結果が観測されるためである。予測より配送が遅れた、交通パターンが違った、制約が現実に合わなかった場合、仮説、情報モデル、意思決定モデルを更新する必要がある。
\end{enumerate}

\end{multicols}

\clearpage

\SlideHeader{Lecture 1: Overview}{018}{アジェンダ}

\SlideImage{18}{../Materials/2026/3DM_01_Overview_SS26.pdf}

\vspace{1.5mm}

\begin{multicols}{2}

\footnotesize

\raggedcolumns

\NoteSection{日本語訳}

\begin{itemize}
\item 動機
\item Operations Research と Business Analytics の復習
\item モデリングと講義全体の見取り図
\end{itemize}

\end{multicols}

\clearpage

\SlideHeader{Lecture 1: Overview}{019}{例: Braunschweigの小包配送}

\SlideImage{19}{../Materials/2026/3DM_01_Overview_SS26.pdf}

\vspace{1.5mm}

\begin{multicols}{2}

\footnotesize

\raggedcolumns

\NoteSection{日本語訳}

\begin{itemize}
\item トラックはデポにいる。
\item 顧客と住所がある。
\item 道路ネットワークがある。
\item 全顧客を訪問するように、市内でのドライバーのナビゲーションを探す。
\item 目標は迅速な配送である。
\item どうすれば目標を達成できるか。
\end{itemize}

\NoteSection{解説}
TSP は、顧客をノード、移動可能性をアーク、移動時間または距離を重みとして表し、全顧客を一度ずつ訪問して出発点に戻る経路を求める問題である。\par

情報モデルは、道路ネットワークや旅行速度を、顧客間の旅行時間行列へ変換する部分である。意思決定モデルは、どのアークを使うかを二値変数で表し、各顧客の入次数・出次数、部分巡回除去、目的関数を定式化する部分である。\par

静的TSPでは旅行時間は定数と仮定する。しかし都市配送では出発時刻により旅行時間が変わるため、この仮説が実装後の遅延や残業として破綻することがある。\par

\NoteSection{試験対策ポイント}
TSPの情報モデルと意思決定モデルを分けて説明し、静的旅行時間仮説が実装後に崩れる理由を説明できること。\par

\NoteSection{過去問関連}
\ExamItem{WS22/23 Aufgabe 1 (7点)}{TSP の数理最適化モデルと、時間次元を持つTSPで何が変わるかを問う問題。}{標準TSPでは、顧客ノード、アーク、旅行時間または距離、二値変数 x\_ij、各顧客を一度訪問する制約、部分巡回除去制約、総旅行時間最小化を定義する。時間依存TSPでは旅行時間が定数 d\_ij ではなく出発時刻 t に依存する d\_ij(t) となり、あるアーク選択が後続の到着時刻と後続アークの旅行時間を変えるため、静的モデルより決定変数・制約・解法が複雑になる。}
\ExamItem{SS24 Aufgabe 1b (4点)}{旅行時間の情報モデルの例と、常に最も詳細なモデルを選ばない理由を問う問題。}{例として、単一平均旅行時間行列、時間帯別行列、区間旅行時間、確率分布モデルがある。詳細なモデルほど現実に近い可能性はあるが、データ量、信頼性、計算時間、解釈可能性、意思決定モデルの複雑性が悪化する。したがって目的に十分な粒度を選ぶ。}
\ExamItem{WS23/24 Aufgabe 2 (5点)}{Real-World Problem と Decision Model の関係、および Information Model を介した流れを説明する問題。}{現実世界から直接最適化モデルを作るのではなく、まず観測データを集約して意思決定に必要な情報表現に落とす。Information Model は現実を圧縮した表現で、Decision Model はその表現を使って決定空間、制約、目的関数を定義する。解は Implementation により現実の行動へ変換され、結果観測から仮説・モデルが更新される。}
\ExamItem{SS24 Aufgabe 1a (8点)}{Information Model, Decision Model, Aggregation, Disaggregation/Implementation の機能を説明する問題。}{Aggregation は現実データをモデルの次元へ変換する。Information Model は意思決定に必要な状態・属性・パラメータを整理する。Decision Model は可能な決定、目的、制約、評価方法を定義する。Disaggregation/Implementation はモデル解を現実の作業指示やルートへ戻す。重要なのは、これらが一方向で終わらず、実装後の観測で更新されるループである点である。}
\ExamItem{WS22/23 Aufgabe 3 (6点)}{Business Analytics プロセスの中で OR 的なモデル化・解法がどこにあるかを説明する問題。}{情報モデルはデータ側、意思決定モデルはOR側の橋渡しである。ORは、現実の問題構造を抽象化し、目的関数・制約・決定変数に変換し、解を評価可能にする。}

\NoteSection{確認問題と解答}
\begin{enumerate}
\item \textbf{TSPの情報モデルと意思決定モデルを分けて説明せよ。}\\
情報モデルは顧客、デポ、顧客間旅行時間行列である。意思決定モデルは、アーク選択変数 x\_ij、各顧客を一度訪問する制約、部分巡回除去制約、総旅行時間最小化の目的関数からなる。
\item \textbf{静的TSPで『旅行時間は一定』という仮説がなぜ問題になるか。}\\
都市交通では朝夕のラッシュ、場所ごとの混雑、曜日差により旅行時間が変わる。旅行時間を一定とすると、モデル上は実行可能で短いルートでも、現実では時間超過や顧客不在が起きる可能性がある。
\end{enumerate}

\end{multicols}

\clearpage

\SlideHeader{Lecture 1: Overview}{020}{モデリング: 巡回セールスマン問題}

\SlideImage{20}{../Materials/2026/3DM_01_Overview_SS26.pdf}

\vspace{1.5mm}

\begin{multicols}{2}

\footnotesize

\raggedcolumns

\NoteSection{日本語訳}

\begin{itemize}
\item Part 1
\begin{itemize}
\item 道路ネットワークを顧客間のアークへ変換する。
\item 走行速度をアークごとの旅行時間へ変換する。
\item 重み付きグラフまたは旅行時間行列を得る。
\end{itemize}
\item Part 2
\begin{itemize}
\item ナビゲーションを、使用アークに関する決定変数へ変換する。
\item 実行可能なナビゲーションを保証する制約を追加する。
\item Part 1に基づいて目的関数を定義する。
\end{itemize}
\end{itemize}

\NoteSection{解説}
TSP は、顧客をノード、移動可能性をアーク、移動時間または距離を重みとして表し、全顧客を一度ずつ訪問して出発点に戻る経路を求める問題である。\par

情報モデルは、道路ネットワークや旅行速度を、顧客間の旅行時間行列へ変換する部分である。意思決定モデルは、どのアークを使うかを二値変数で表し、各顧客の入次数・出次数、部分巡回除去、目的関数を定式化する部分である。\par

静的TSPでは旅行時間は定数と仮定する。しかし都市配送では出発時刻により旅行時間が変わるため、この仮説が実装後の遅延や残業として破綻することがある。\par

\NoteSection{試験対策ポイント}
TSPの情報モデルと意思決定モデルを分けて説明し、静的旅行時間仮説が実装後に崩れる理由を説明できること。\par

\NoteSection{過去問関連}
\ExamItem{WS22/23 Aufgabe 1 (7点)}{TSP の数理最適化モデルと、時間次元を持つTSPで何が変わるかを問う問題。}{標準TSPでは、顧客ノード、アーク、旅行時間または距離、二値変数 x\_ij、各顧客を一度訪問する制約、部分巡回除去制約、総旅行時間最小化を定義する。時間依存TSPでは旅行時間が定数 d\_ij ではなく出発時刻 t に依存する d\_ij(t) となり、あるアーク選択が後続の到着時刻と後続アークの旅行時間を変えるため、静的モデルより決定変数・制約・解法が複雑になる。}
\ExamItem{SS24 Aufgabe 1b (4点)}{旅行時間の情報モデルの例と、常に最も詳細なモデルを選ばない理由を問う問題。}{例として、単一平均旅行時間行列、時間帯別行列、区間旅行時間、確率分布モデルがある。詳細なモデルほど現実に近い可能性はあるが、データ量、信頼性、計算時間、解釈可能性、意思決定モデルの複雑性が悪化する。したがって目的に十分な粒度を選ぶ。}
\ExamItem{WS23/24 Aufgabe 2 (5点)}{Real-World Problem と Decision Model の関係、および Information Model を介した流れを説明する問題。}{現実世界から直接最適化モデルを作るのではなく、まず観測データを集約して意思決定に必要な情報表現に落とす。Information Model は現実を圧縮した表現で、Decision Model はその表現を使って決定空間、制約、目的関数を定義する。解は Implementation により現実の行動へ変換され、結果観測から仮説・モデルが更新される。}
\ExamItem{SS24 Aufgabe 1a (8点)}{Information Model, Decision Model, Aggregation, Disaggregation/Implementation の機能を説明する問題。}{Aggregation は現実データをモデルの次元へ変換する。Information Model は意思決定に必要な状態・属性・パラメータを整理する。Decision Model は可能な決定、目的、制約、評価方法を定義する。Disaggregation/Implementation はモデル解を現実の作業指示やルートへ戻す。重要なのは、これらが一方向で終わらず、実装後の観測で更新されるループである点である。}
\ExamItem{WS22/23 Aufgabe 3 (6点)}{Business Analytics プロセスの中で OR 的なモデル化・解法がどこにあるかを説明する問題。}{情報モデルはデータ側、意思決定モデルはOR側の橋渡しである。ORは、現実の問題構造を抽象化し、目的関数・制約・決定変数に変換し、解を評価可能にする。}

\NoteSection{確認問題と解答}
\begin{enumerate}
\item \textbf{TSPの情報モデルと意思決定モデルを分けて説明せよ。}\\
情報モデルは顧客、デポ、顧客間旅行時間行列である。意思決定モデルは、アーク選択変数 x\_ij、各顧客を一度訪問する制約、部分巡回除去制約、総旅行時間最小化の目的関数からなる。
\item \textbf{静的TSPで『旅行時間は一定』という仮説がなぜ問題になるか。}\\
都市交通では朝夕のラッシュ、場所ごとの混雑、曜日差により旅行時間が変わる。旅行時間を一定とすると、モデル上は実行可能で短いルートでも、現実では時間超過や顧客不在が起きる可能性がある。
\end{enumerate}

\end{multicols}

\clearpage

\SlideHeader{Lecture 1: Overview}{021}{解を見つける}

\SlideImage{21}{../Materials/2026/3DM_01_Overview_SS26.pdf}

\vspace{1.5mm}

\begin{multicols}{2}

\footnotesize

\raggedcolumns

\NoteSection{日本語訳}

\begin{itemize}
\item モデルは解空間を張り、解を評価する。
\item 候補解の求め方
\begin{itemize}
\item グラフベース手法、例: insertion procedures
\item 混合整数計画、例: branch and bound
\end{itemize}
\item 最後のステップ: 実務への実装
\begin{itemize}
\item 数字を顧客へ変換する。
\item 選択されたアークを道路セグメントの集合へ対応させる。
\item ドライバーにナビゲーションを提供する。
\end{itemize}
\end{itemize}

\NoteSection{解説}
TSP は、顧客をノード、移動可能性をアーク、移動時間または距離を重みとして表し、全顧客を一度ずつ訪問して出発点に戻る経路を求める問題である。\par

情報モデルは、道路ネットワークや旅行速度を、顧客間の旅行時間行列へ変換する部分である。意思決定モデルは、どのアークを使うかを二値変数で表し、各顧客の入次数・出次数、部分巡回除去、目的関数を定式化する部分である。\par

静的TSPでは旅行時間は定数と仮定する。しかし都市配送では出発時刻により旅行時間が変わるため、この仮説が実装後の遅延や残業として破綻することがある。\par

\NoteSection{試験対策ポイント}
TSPの情報モデルと意思決定モデルを分けて説明し、静的旅行時間仮説が実装後に崩れる理由を説明できること。\par

\NoteSection{過去問関連}
\ExamItem{WS22/23 Aufgabe 1 (7点)}{TSP の数理最適化モデルと、時間次元を持つTSPで何が変わるかを問う問題。}{標準TSPでは、顧客ノード、アーク、旅行時間または距離、二値変数 x\_ij、各顧客を一度訪問する制約、部分巡回除去制約、総旅行時間最小化を定義する。時間依存TSPでは旅行時間が定数 d\_ij ではなく出発時刻 t に依存する d\_ij(t) となり、あるアーク選択が後続の到着時刻と後続アークの旅行時間を変えるため、静的モデルより決定変数・制約・解法が複雑になる。}
\ExamItem{SS24 Aufgabe 1b (4点)}{旅行時間の情報モデルの例と、常に最も詳細なモデルを選ばない理由を問う問題。}{例として、単一平均旅行時間行列、時間帯別行列、区間旅行時間、確率分布モデルがある。詳細なモデルほど現実に近い可能性はあるが、データ量、信頼性、計算時間、解釈可能性、意思決定モデルの複雑性が悪化する。したがって目的に十分な粒度を選ぶ。}
\ExamItem{WS23/24 Aufgabe 2 (5点)}{Real-World Problem と Decision Model の関係、および Information Model を介した流れを説明する問題。}{現実世界から直接最適化モデルを作るのではなく、まず観測データを集約して意思決定に必要な情報表現に落とす。Information Model は現実を圧縮した表現で、Decision Model はその表現を使って決定空間、制約、目的関数を定義する。解は Implementation により現実の行動へ変換され、結果観測から仮説・モデルが更新される。}
\ExamItem{SS24 Aufgabe 1a (8点)}{Information Model, Decision Model, Aggregation, Disaggregation/Implementation の機能を説明する問題。}{Aggregation は現実データをモデルの次元へ変換する。Information Model は意思決定に必要な状態・属性・パラメータを整理する。Decision Model は可能な決定、目的、制約、評価方法を定義する。Disaggregation/Implementation はモデル解を現実の作業指示やルートへ戻す。重要なのは、これらが一方向で終わらず、実装後の観測で更新されるループである点である。}
\ExamItem{WS22/23 Aufgabe 3 (6点)}{Business Analytics プロセスの中で OR 的なモデル化・解法がどこにあるかを説明する問題。}{情報モデルはデータ側、意思決定モデルはOR側の橋渡しである。ORは、現実の問題構造を抽象化し、目的関数・制約・決定変数に変換し、解を評価可能にする。}

\NoteSection{確認問題と解答}
\begin{enumerate}
\item \textbf{TSPの情報モデルと意思決定モデルを分けて説明せよ。}\\
情報モデルは顧客、デポ、顧客間旅行時間行列である。意思決定モデルは、アーク選択変数 x\_ij、各顧客を一度訪問する制約、部分巡回除去制約、総旅行時間最小化の目的関数からなる。
\item \textbf{静的TSPで『旅行時間は一定』という仮説がなぜ問題になるか。}\\
都市交通では朝夕のラッシュ、場所ごとの混雑、曜日差により旅行時間が変わる。旅行時間を一定とすると、モデル上は実行可能で短いルートでも、現実では時間超過や顧客不在が起きる可能性がある。
\end{enumerate}

\end{multicols}

\clearpage

\SlideHeader{Lecture 1: Overview}{022}{Operations Researchの三段階}

\SlideImage{22}{../Materials/2026/3DM_01_Overview_SS26.pdf}

\vspace{1.5mm}

\begin{multicols}{2}

\footnotesize

\raggedcolumns

\NoteSection{日本語訳}

\begin{itemize}
\item 問題定義
\item モデリング
\item 解法
\item 情報と意思決定
\item 実装を含む。
\end{itemize}

\NoteSection{解説}
IDA は、観測されたデータの appearance から、背後にある system structure を推測する方向である。データのクリーニング、補完、検証、データマイニング、パターン抽出が含まれる。\par

PMT/OR は、問題構造から意思決定モデルを作り、解を計算して現実に実装する方向である。目的、決定空間、制約、ルーティング・スケジューリング・割当のようなモデル構造が重要になる。\par

統一的に見ると、Business Analytics はデータから情報モデルを作る流れと、問題構造から意思決定モデルを作る流れを接続する。DDDMで重要なのは、どちらか一方ではなく、情報モデルと意思決定モデルが互いに依存する点である。\par

\NoteSection{試験対策ポイント}
IDAはデータから構造へ、OR/PMTは構造から解へ向かう。この二つがBusiness Analyticsで接続されることを図なしで説明できること。\par

\NoteSection{過去問関連}
\ExamItem{WS22/23 Aufgabe 2 (6点)}{Business Analytics を構成する分野と、それらの交差領域を説明する問題。}{Operations Research は意思決定モデル、制約、目的関数、最適化を担う。Statistics はデータの不確実性、分布、推定、予測を扱う。Computer Science/Information Systems はデータ管理、アルゴリズム、実装、情報システムを担う。Business Analytics はこれらを組み合わせ、データから意思決定可能な知見と実行可能な解を作る分野である。}
\ExamItem{WS22/23 Aufgabe 3 (6点)}{Business Analytics の適用プロセスを挙げ、Operations Research に属する工程を示す問題。}{現実問題を観察し、データを取得・集約して情報モデルを作り、そこから意思決定モデルを定義し、解法で解を求め、実装して現実に戻す。ORに特に属するのは、目的関数・制約・決定変数を持つ意思決定モデルの定式化、解空間の探索、最適化・ヒューリスティックによる解の計算である。}
\ExamItem{WS23/24 Aufgabe 1 (6点)}{Business Analytics の三つの構成分野と交差領域を説明する問題。}{三分野を単に列挙するだけでなく、データ分析だけでは意思決定にならず、最適化だけでも現実データを反映できない点を述べる。交差領域では、データからモデルを作り、モデルから現実に実装可能な意思決定を作ることが中心になる。}
\ExamItem{WS23/24 Aufgabe 3 (10点)}{配送車両の業務的ツアープランニングに Business Analytics の構成要素を対応付ける問題。}{Real-world は配送現場、顧客、道路、車両である。Aggregation は住所・道路・走行時間観測をモデル変数へ変換すること。Information model は顧客集合、距離/旅行時間行列、時間帯情報である。Decision model は訪問順序・経路・目的関数・制約を持つVRP/TSPである。Implementation は得られたツアーをドライバーのルート案内として実行することである。}
\ExamItem{WS23/24 Aufgabe 2 (5点)}{Real-World Problem と Decision Model の関係、および Information Model を介した流れを説明する問題。}{現実世界から直接最適化モデルを作るのではなく、まず観測データを集約して意思決定に必要な情報表現に落とす。Information Model は現実を圧縮した表現で、Decision Model はその表現を使って決定空間、制約、目的関数を定義する。解は Implementation により現実の行動へ変換され、結果観測から仮説・モデルが更新される。}
\ExamItem{SS24 Aufgabe 1a (8点)}{Information Model, Decision Model, Aggregation, Disaggregation/Implementation の機能を説明する問題。}{Aggregation は現実データをモデルの次元へ変換する。Information Model は意思決定に必要な状態・属性・パラメータを整理する。Decision Model は可能な決定、目的、制約、評価方法を定義する。Disaggregation/Implementation はモデル解を現実の作業指示やルートへ戻す。重要なのは、これらが一方向で終わらず、実装後の観測で更新されるループである点である。}
\ExamItem{WS22/23 Aufgabe 3 (6点)}{Business Analytics プロセスの中で OR 的なモデル化・解法がどこにあるかを説明する問題。}{情報モデルはデータ側、意思決定モデルはOR側の橋渡しである。ORは、現実の問題構造を抽象化し、目的関数・制約・決定変数に変換し、解を評価可能にする。}

\NoteSection{確認問題と解答}
\begin{enumerate}
\item \textbf{IDA と OR/PMT の方向性の違いを説明せよ。}\\
IDAは観測データから構造を見つける。つまり appearance から structure へ向かう。OR/PMTは問題構造を仮説として選び、意思決定モデルと解を作る。つまり structure から solution/appearance へ向かう。Business Analyticsでは両者を循環的に使う。
\item \textbf{なぜデータはそのまま構造を示さないのか。}\\
データは現実の一部の観測であり、不完全、誤り、例外、ノイズを含む。例えばGoogle Mapsが渋滞を示しても、その原因や構造までは直接示さない。したがってデータ処理と仮説に基づくモデル化が必要になる。
\end{enumerate}

\end{multicols}

\clearpage

\SlideHeader{Lecture 1: Overview}{023}{アジェンダ}

\SlideImage{23}{../Materials/2026/3DM_01_Overview_SS26.pdf}

\vspace{1.5mm}

\begin{multicols}{2}

\footnotesize

\raggedcolumns

\NoteSection{日本語訳}

\begin{itemize}
\item 動機
\item Operations Research と Business Analytics の復習
\item モデリングと講義全体の見取り図
\end{itemize}

\end{multicols}

\clearpage

\SlideHeader{Lecture 1: Overview}{024}{モデリング}

\SlideImage{24}{../Materials/2026/3DM_01_Overview_SS26.pdf}

\vspace{1.5mm}

\begin{multicols}{2}

\footnotesize

\raggedcolumns

\NoteSection{日本語訳}

\begin{itemize}
\item 異なる問題が同じモデルを共有する場合がある。例: 小包配送と技術者サービス。
\item 同じ問題を異なる方法でモデル化することもできる。
\item 情報モデルと意思決定モデルの間には接続がある。
\end{itemize}

\NoteSection{解説}
Real-World は、配送、交通、顧客、車両、作業員など、実際に行動が行われる対象である。ここでは全ての詳細をモデルに入れることはできないため、意思決定に必要な情報だけを選ぶ。\par

Aggregation は、観測データをモデルで使える次元へ変換する処理である。例えばGPS位置列を道路セグメント別の平均速度へ変換すること、住所を顧客ノードへ変換することが該当する。\par

Information Model は、現実を圧縮した情報表現である。旅行時間行列、時間帯別速度パターン、顧客集合、需要分布などが含まれる。\par

Decision Model は、可能な決定、目的関数、制約、評価方法を定義する。TSP/VRPであれば、訪問順序、利用アーク、各顧客を一度訪問する制約、総旅行時間最小化が含まれる。\par

Implementation は、モデルの解を現実の行動へ戻す処理である。例えば『1-3-2』というノード順序を、実際の住所と道路ナビゲーションへ変換する。\par

\NoteSection{試験対策ポイント}
Real-World, Aggregation, Information Model, Decision Model, Implementation の役割とデータ・解の流れを、配送例に対応付けて説明できること。\par

\NoteSection{過去問関連}
\ExamItem{WS23/24 Aufgabe 2 (5点)}{Real-World Problem と Decision Model の関係、および Information Model を介した流れを説明する問題。}{現実世界から直接最適化モデルを作るのではなく、まず観測データを集約して意思決定に必要な情報表現に落とす。Information Model は現実を圧縮した表現で、Decision Model はその表現を使って決定空間、制約、目的関数を定義する。解は Implementation により現実の行動へ変換され、結果観測から仮説・モデルが更新される。}
\ExamItem{SS24 Aufgabe 1a (8点)}{Information Model, Decision Model, Aggregation, Disaggregation/Implementation の機能を説明する問題。}{Aggregation は現実データをモデルの次元へ変換する。Information Model は意思決定に必要な状態・属性・パラメータを整理する。Decision Model は可能な決定、目的、制約、評価方法を定義する。Disaggregation/Implementation はモデル解を現実の作業指示やルートへ戻す。重要なのは、これらが一方向で終わらず、実装後の観測で更新されるループである点である。}
\ExamItem{WS22/23 Aufgabe 3 (6点)}{Business Analytics プロセスの中で OR 的なモデル化・解法がどこにあるかを説明する問題。}{情報モデルはデータ側、意思決定モデルはOR側の橋渡しである。ORは、現実の問題構造を抽象化し、目的関数・制約・決定変数に変換し、解を評価可能にする。}
\ExamItem{WS22/23 Aufgabe 2 (6点)}{Business Analytics を構成する分野と、それらの交差領域を説明する問題。}{Operations Research は意思決定モデル、制約、目的関数、最適化を担う。Statistics はデータの不確実性、分布、推定、予測を扱う。Computer Science/Information Systems はデータ管理、アルゴリズム、実装、情報システムを担う。Business Analytics はこれらを組み合わせ、データから意思決定可能な知見と実行可能な解を作る分野である。}
\ExamItem{WS22/23 Aufgabe 3 (6点)}{Business Analytics の適用プロセスを挙げ、Operations Research に属する工程を示す問題。}{現実問題を観察し、データを取得・集約して情報モデルを作り、そこから意思決定モデルを定義し、解法で解を求め、実装して現実に戻す。ORに特に属するのは、目的関数・制約・決定変数を持つ意思決定モデルの定式化、解空間の探索、最適化・ヒューリスティックによる解の計算である。}
\ExamItem{WS23/24 Aufgabe 1 (6点)}{Business Analytics の三つの構成分野と交差領域を説明する問題。}{三分野を単に列挙するだけでなく、データ分析だけでは意思決定にならず、最適化だけでも現実データを反映できない点を述べる。交差領域では、データからモデルを作り、モデルから現実に実装可能な意思決定を作ることが中心になる。}
\ExamItem{WS23/24 Aufgabe 3 (10点)}{配送車両の業務的ツアープランニングに Business Analytics の構成要素を対応付ける問題。}{Real-world は配送現場、顧客、道路、車両である。Aggregation は住所・道路・走行時間観測をモデル変数へ変換すること。Information model は顧客集合、距離/旅行時間行列、時間帯情報である。Decision model は訪問順序・経路・目的関数・制約を持つVRP/TSPである。Implementation は得られたツアーをドライバーのルート案内として実行することである。}

\NoteSection{確認問題と解答}
\begin{enumerate}
\item \textbf{Information Model と Decision Model の違いを説明せよ。}\\
Information Model は、意思決定に必要な現実情報の表現である。例として顧客、道路、旅行時間行列、時間帯別速度がある。Decision Model は、その情報を用いて決定を選ぶための目的関数、制約、決定変数の体系である。情報モデルが『何を知っているか』を表し、意思決定モデルが『その情報で何をどう決めるか』を表す。
\item \textbf{Aggregation と Implementation の役割を説明せよ。}\\
Aggregation は現実データを抽象化してモデル入力に変換する処理であり、Implementation はモデル解を実際の行動へ具体化する処理である。前者は現実からモデルへの橋、後者はモデルから現実への橋である。
\item \textbf{なぜこのループは一回で終わらないのか。}\\
実装後に現実の結果が観測されるためである。予測より配送が遅れた、交通パターンが違った、制約が現実に合わなかった場合、仮説、情報モデル、意思決定モデルを更新する必要がある。
\end{enumerate}

\end{multicols}

\clearpage

\SlideHeader{Lecture 1: Overview}{025}{例I: 代替的な情報モデル}

\SlideImage{25}{../Materials/2026/3DM_01_Overview_SS26.pdf}

\vspace{1.5mm}

\begin{multicols}{2}

\footnotesize

\raggedcolumns

\NoteSection{日本語訳}

\begin{itemize}
\item 旅行時間の日内変動、つまり時間依存旅行時間を観測する。
\item 7-9時、9-15時、15-18時のような時間帯を考える。
\item 情報モデルは複数の旅行時間行列の集合を含む。
\item 情報生成はより複雑になる。
\item 代替的な意思決定モデルが必要になる。
\end{itemize}

\NoteSection{解説}
時間依存旅行時間では、顧客 i から j への旅行時間を定数ではなく出発時刻 t の関数として扱う。記号的には d\_ij(t) または tau\_ij(t) と考える。\par

同じ二地点間でも、出発時刻が異なれば最短経路も旅行時間も変わる。したがって、訪問順序の変更は単にアークの足し算を変えるだけでなく、後続の到着時刻と全体の旅行時間を変える。\par

このため、情報モデルは時間帯別行列や時間依存リンク旅行時間を必要とし、意思決定モデルは時間の次元を含む変数・制約を必要とする。\par

\NoteSection{試験対策ポイント}
旅行時間を d\_ij(t) として捉え、出発時刻が経路・順序・後続旅行時間に与える影響を説明できること。\par

\NoteSection{過去問関連}
\ExamItem{WS22/23 Aufgabe 1 (7点)}{TSP の数理最適化モデルと、時間次元を持つTSPで何が変わるかを問う問題。}{標準TSPでは、顧客ノード、アーク、旅行時間または距離、二値変数 x\_ij、各顧客を一度訪問する制約、部分巡回除去制約、総旅行時間最小化を定義する。時間依存TSPでは旅行時間が定数 d\_ij ではなく出発時刻 t に依存する d\_ij(t) となり、あるアーク選択が後続の到着時刻と後続アークの旅行時間を変えるため、静的モデルより決定変数・制約・解法が複雑になる。}
\ExamItem{SS24 Aufgabe 1b (4点)}{旅行時間の情報モデルの例と、常に最も詳細なモデルを選ばない理由を問う問題。}{例として、単一平均旅行時間行列、時間帯別行列、区間旅行時間、確率分布モデルがある。詳細なモデルほど現実に近い可能性はあるが、データ量、信頼性、計算時間、解釈可能性、意思決定モデルの複雑性が悪化する。したがって目的に十分な粒度を選ぶ。}
\ExamItem{WS22/23 Aufgabe 4 (8点)}{過去の交通観測から時間帯別平均旅行時間を構築する手順を説明する問題。}{まずFCDなどから位置・時刻・速度を収集する。位置を道路セグメントにマッチングし、セグメントごとの観測速度を得る。外れ値やエラーを除去し、時間帯 w とセグメント l ごとに観測を集約する。平均速度または旅行時間を計算し、観測不足の部分は類似セグメントや時間帯で一般化する。最後に時間帯別の旅行時間行列を作り、TSP/VRPの情報モデルとして使う。}
\ExamItem{SS24 Aufgabe 1b (4点)}{旅行時間情報モデルの選択と複雑性のトレードオフを問う問題。}{平均値モデルは簡単だがラッシュアワーを無視しやすい。時間帯別モデルは実際の交通変動を反映するが、観測数が少ないセルでは信頼性が低くなる。したがってデータ量と意思決定への効果を見ながら集約・一般化する必要がある。}

\NoteSection{確認問題と解答}
\begin{enumerate}
\item \textbf{時間依存旅行時間モデルでは、なぜ一つの旅行時間行列では不十分か。}\\
一つの行列では全時間帯で同じ旅行時間を仮定することになる。実際には道路セグメントごとに混雑時間が異なり、10時と17時では同じ顧客間の最短経路や所要時間が異なる。そのため時間帯別または連続時刻の情報モデルが必要になる。
\item \textbf{時間依存性がヒューリスティックを難しくする理由を説明せよ。}\\
ヒューリスティックの局所的な選択が、後続の出発時刻を変え、後続の旅行時間関数の評価値を変えるからである。静的な追加距離だけでは全体影響を評価できない。
\end{enumerate}

\end{multicols}

\clearpage

\SlideHeader{Lecture 1: Overview}{026}{例II: 代替的な意思決定モデル}

\SlideImage{26}{../Materials/2026/3DM_01_Overview_SS26.pdf}

\vspace{1.5mm}

\begin{multicols}{2}

\footnotesize

\raggedcolumns

\NoteSection{日本語訳}

\begin{itemize}
\item 解が実行可能に見えても、残業が頻繁に発生する。
\item 旅行時間に不確実性がある。
\item 意思決定モデルを変更する。これは stochasticity に関係する。
\item 代替的でより詳細な情報モデルが必要になる。
\end{itemize}

\NoteSection{解説}
不確実性とは、意思決定時点では正確に分からないが、結果や実行可能性に影響する情報である。例として需要、株価、交通、天候、サービス時間、ドライバー利用可能性などがある。\par

不確実性は頻度、変動幅、破壊力、予測可能性によって性質が異なる。小さく頻繁な変動と、稀だが大きな事故・災害では、モデル化と意思決定方法が異なる。\par

DDDMで重要なのは、不確実性を情報モデルにどう表すか、そして意思決定モデルでどう評価するかである。\par

\NoteSection{試験対策ポイント}
不確実性の種類と、計画失敗・遅延・追加費用などの意思決定上の影響を具体例で説明できること。\par

\NoteSection{過去問関連}
\ExamItem{WS22/23 Aufgabe 5 (6点)}{旅行時間の不確実性、確率的/準確率的モデル、ロバスト解の違いを問う問題。}{確率的モデルでは確率分布があり、期待値・分散・サンプリングを使って判断する。準確率的モデルでは確率が分からず、シナリオ、区間、最悪ケース、Hurwicz、regret などで不確実性を扱う。ロバスト解は、単一の平均ケースに最適化するのではなく、複数のあり得る状況で大きく破綻しにくい解である。}
\ExamItem{WS24/25 Aufgabe 3 (12点)}{stochastic と quasi-stochastic の違い、および各手法を説明する問題。}{stochastic は実現値の確率分布または確率が既知で、期待値最小化、Monte Carlo sampling、期待値-分散基準などを使える。quasi-stochastic は確率が未知で、シナリオ集合、区間、Hurwicz基準、minmax regret、最悪ケース最適化などを使う。判断の基準は、データ量、確率分布の妥当性、リスク回避度、解釈可能性である。}
\ExamItem{SS24 Aufgabe 2a (3点)}{Same-Day配送で考えられる不確実性のカテゴリを挙げる問題。}{需要不確実性として注文発生・量・場所、リソース不確実性としてドライバーの利用可能性・スキル・車両/バッテリー、環境不確実性として交通、駐車、天候、道路障害を挙げる。}

\NoteSection{確認問題と解答}
\begin{enumerate}
\item \textbf{物流・モビリティにおける不確実性を三分類で挙げよ。}\\
顧客側の不確実性として注文発生、需要量、サービス時間、在宅状況がある。ドライバー・リソース側の不確実性として可用性、行動、航続距離、スキルがある。環境側の不確実性として交通、駐車、道路ネットワーク、天候がある。
\item \textbf{不確実性を考慮しないと何が起こるか。}\\
計画が実行不能になる、遅延が発生する、顧客が不在になり配送失敗になる、追加費用が発生する、後続活動が連鎖的に遅れる。したがって不確実性は単なる誤差ではなく、意思決定品質を大きく左右する要素である。
\end{enumerate}

\end{multicols}

\clearpage

\SlideHeader{Lecture 1: Overview}{027}{例III: 動的性}

\SlideImage{27}{../Materials/2026/3DM_01_Overview_SS26.pdf}

\vspace{1.5mm}

\begin{multicols}{2}

\footnotesize

\raggedcolumns

\NoteSection{日本語訳}

\begin{itemize}
\item 情報モデルが時間とともに変化することを観測する。
\item 変化に反応できる。
\item 確率的・動的モデルが必要になる。
\item 将来の先読みが重要になる。
\item Approximate Dynamic Programming へつながる。
\end{itemize}

\NoteSection{解説}
DDDM は、データを単に可視化・分析するだけでなく、そのデータを使って具体的な行動を選ぶための枠組みである。したがって中心は『どの情報を使い、どのモデルで、どの決定を選び、現実にどう実装するか』である。\par

stochasticity は、需要、交通、サービス時間などの情報が不確実で、時間とともに新しく明らかになることを指す。dynamism は、新しい情報が出た後に計画を更新しながら意思決定を続ける性質を指す。\par

\NoteSection{試験対策ポイント}
DDDMは、データ分析、意思決定モデル、現実への実装を結ぶ講義である。stochasticity と dynamism を必ず区別して説明できるようにする。\par

\NoteSection{確認問題と解答}
\begin{enumerate}
\item \textbf{DDDM は通常のデータ分析と何が違うか。}\\
通常のデータ分析は、過去データの説明や予測で終わることがある。DDDMでは、その情報を使って実際の決定、例えば配送順序、車両割当、注文受諾、在庫補充などを選ぶ。したがって、情報モデルと意思決定モデル、さらに現実への実装までを一体として考える。
\item \textbf{stochasticity と dynamism を区別して説明せよ。}\\
stochasticity は将来情報や実現値が不確実であること、dynamism は時間の経過とともに状態が変わり、そのたびに決定を更新する必要があることである。例えば、注文が確率的に発生することは stochasticity、注文発生後に配送計画を組み替えることは dynamism である。
\end{enumerate}

\end{multicols}

\clearpage

\SlideHeader{Lecture 1: Overview}{028}{次回予告}

\SlideImage{28}{../Materials/2026/3DM_01_Overview_SS26.pdf}

\vspace{1.5mm}

\begin{multicols}{2}

\footnotesize

\raggedcolumns

\NoteSection{日本語訳}

\begin{itemize}
\item Business Analytics
\item 情報モデリングと意思決定モデリング
\item ケーススタディ: 時間依存旅行時間を持つ配送ルーティング
\end{itemize}

\end{multicols}
\end{document}
